\documentclass[11pt]{article}
% ===========================================================================
%  Shared preamble for the AP Calculus AB exam-preparation worksheets.
%  Input right after \documentclass. Files beginning with "_" are not compiled
%  on their own. Navy and gold, to match the course's lesson pages.
% ===========================================================================
\usepackage[letterpaper,top=2.3cm,bottom=2.2cm,left=1.7cm,right=1.7cm,headheight=28pt,footskip=30pt]{geometry}
\usepackage{amsmath,amssymb}
\usepackage[T1]{fontenc}
\usepackage[utf8]{inputenc}
\usepackage{lmodern}
\usepackage{xcolor}
\usepackage[most]{tcolorbox}
\usepackage{enumitem}
\usepackage{fancyhdr}
\usepackage{array}
\usepackage{tikz}
\usepackage{pgfplots}
\pgfplotsset{compat=1.18}

\definecolor{navy}{HTML}{161E54}
\definecolor{ink}{HTML}{27357E}
\definecolor{gold}{HTML}{B98A00}
\definecolor{goldbg}{HTML}{FFFDF3}
\definecolor{inkbg}{HTML}{F6F7FD}
\definecolor{rule}{HTML}{C9CFF0}
\definecolor{okgreen}{HTML}{1B7A44}
\definecolor{okbg}{HTML}{F1F8F3}

\pagestyle{fancy}
\fancyhf{}
\fancyhead[L]{\footnotesize Name: \underline{\hspace{4cm}}}
\fancyhead[C]{\textbf{\color{navy}Integration Academy}\\[-2pt]{\scriptsize Calculus AB \textperiodcentered{} AP\textsuperscript{\textregistered} Exam Preparation}}
\fancyfoot[L]{\tiny\color{black!55}\parbox[t]{0.78\textwidth}{AP\textsuperscript{\textregistered} is a registered trademark of the College Board, which was not involved in the production of, and does not endorse, this material. \textcopyright{} Integration Academy.}}
\fancyfoot[R]{\scriptsize\thepage}
\renewcommand{\headrulewidth}{1.2pt}
\renewcommand{\headrule}{\hbox to\headwidth{\color{gold}\leaders\hrule height \headrulewidth\hfill}}
\renewcommand{\footrulewidth}{0pt}
\setlength{\parindent}{0pt}
\setlength{\parskip}{4pt}

% Title block: \wstitle{code}{title}{unit line}
\newcommand{\wstitle}[3]{\begin{center}{\small\bfseries\color{gold}WORKSHEET #1}\\[3pt]{\LARGE\bfseries\color{navy}#2}\\[3pt]{\small\color{black!60}#3}\end{center}\vspace{2pt}}
% Section heading
\newcommand{\sech}[1]{\par\vspace{10pt}{\large\bfseries\color{navy}#1}\par\nobreak\vspace{2pt}{\color{rule}\hrule height0.8pt}\vspace{7pt}}
% A level tag printed beside a group of questions
\newcommand{\level}[1]{\par\vspace{6pt}{\footnotesize\bfseries\color{gold}\MakeUppercase{#1}}\par\nobreak\vspace{2pt}}

\newtcolorbox{ideabox}{colframe=ink,colback=inkbg,coltitle=white,colbacktitle=ink,fonttitle=\bfseries\small,title={Key ideas},arc=2pt,boxrule=0.8pt,left=8pt,right=8pt}
\newtcolorbox{example}[1]{colframe=rule,colback=white,coltitle=navy,colbacktitle=inkbg,fonttitle=\bfseries\small,title={#1},arc=2pt,boxrule=0.8pt,borderline west={3pt}{0pt}{ink},left=10pt,right=8pt}
\newtcolorbox{frbox}[1]{colframe=gold,colback=goldbg,coltitle=white,colbacktitle=gold,fonttitle=\bfseries\small,title={#1},arc=2pt,boxrule=0.8pt,breakable,left=8pt,right=8pt}
\newtcolorbox{solbox}{colframe=okgreen,colback=okbg,coltitle=white,colbacktitle=okgreen,fonttitle=\bfseries\small,title={Worked solutions},arc=2pt,boxrule=0.8pt,breakable,left=8pt,right=8pt}
\newtcolorbox{scorebox}[1]{colframe=navy,colback=white,coltitle=white,colbacktitle=navy,fonttitle=\bfseries\small,title={#1},arc=2pt,boxrule=0.8pt,breakable,left=8pt,right=8pt}

\newcommand{\soln}{\par\smallskip\textbf{\color{ink}Solution.}\ }
\newcommand{\result}[1]{\par\smallskip\textbf{#1}}
% blank working space of a given height
\newcommand{\work}[1]{\par\vspace{#1}}
% a point value in the margin of a scoring guide
\newcommand{\pts}[1]{\hfill{\footnotesize\bfseries\color{gold}[#1]}}

% Coordinate grid for graph questions: \begin{apgraph}{xmin}{xmax}{ymin}{ymax} ... \end{apgraph}
\newenvironment{apgraph}[4]{\begin{center}\begin{tikzpicture}\begin{axis}[
  axis lines=middle, xmin=#1, xmax=#2, ymin=#3, ymax=#4,
  grid=both, grid style={black!12}, xtick distance=1, ytick distance=1,
  ticklabel style={font=\scriptsize}, xlabel={$x$}, ylabel={$y$},
  width=10cm, height=7.2cm, samples=80, clip=false,
  every axis plot/.append style={line width=1.1pt, color=ink}]}{\end{axis}\end{tikzpicture}\end{center}}
\newcommand{\closeddot}[2]{\addplot[only marks, mark=*, mark size=2.2pt, color=ink] coordinates {(#1,#2)};}
\newcommand{\opendot}[2]{\addplot[only marks, mark=*, mark size=2.2pt, color=ink, mark options={fill=white}] coordinates {(#1,#2)};}

% Multiple-choice options: four across, or stacked when the options are long.
\newcommand{\opts}[4]{\par\vspace{3pt}\begin{tabular}{@{}p{0.235\linewidth}p{0.235\linewidth}p{0.235\linewidth}p{0.235\linewidth}@{}}(A)\ #1 & (B)\ #2 & (C)\ #3 & (D)\ #4\end{tabular}\par\vspace{9pt}}
\newcommand{\optsv}[4]{\par\vspace{3pt}\begin{tabular}{@{}l@{\ \ }p{0.9\linewidth}@{}}(A) & #1\\(B) & #2\\(C) & #3\\(D) & #4\end{tabular}\par\vspace{9pt}}

\fancyhead[R]{\footnotesize\color{black!60}Worksheet 6.9}
\begin{document}

\wstitle{6.9}{Integration by Substitution}{Unit 6: Integration and Accumulation of Change}

\begin{ideabox}
\begin{itemize}[leftmargin=*,itemsep=2pt,topsep=2pt]
  \item Substitution undoes the chain rule. Choose $u$ to be an \textbf{inside} function whose derivative appears as a factor.
  \item Replace every $x$ and the $dx$: if $u=g(x)$ then $du=g'(x)\,dx$.
  \item A missing constant factor can be supplied: if $du=2x\,dx$, then $x\,dx=\dfrac12du$.
  \item In a definite integral, \textbf{change the limits} to values of $u$, or return to $x$ before evaluating. Never mix the two.
\end{itemize}
\end{ideabox}

\sech{Worked examples}
\begin{example}{Example 1 \textperiodcentered{} A direct fit}
Find $\displaystyle\int3x^2\left(x^3+5\right)^4dx$.\soln $u=x^3+5$, $du=3x^2\,dx$: $\displaystyle\int u^4\,du$.\result{$\dfrac{\left(x^3+5\right)^5}{5}+C$.}
\end{example}

\begin{example}{Example 2 \textperiodcentered{} A linear inside}
Find $\displaystyle\int\cos(4x)\,dx$.\soln $u=4x$, $dx=\dfrac14du$.\result{$\dfrac14\sin(4x)+C$.}
\end{example}

\begin{example}{Example 3 \textperiodcentered{} Changing the limits}
Evaluate $\displaystyle\int_{0}^{3}\dfrac{2x}{x^2+1}\,dx$.\soln $u=x^2+1$, $du=2x\,dx$; when $x=0$, $u=1$; when $x=3$, $u=10$. $\displaystyle\int_{1}^{10}\dfrac1u\,du$.\result{$\ln10$.}
\end{example}

\newpage
\sech{Practice}

\level{Core}
\begin{enumerate}[leftmargin=*,itemsep=2pt]
  \item Find $\displaystyle\int2x\left(x^2+1\right)^3dx$.\work{2.3cm}
  \item Find $\displaystyle\int e^{5x}\,dx$.\work{2.3cm}
  \item Find $\displaystyle\int\sin(3x)\,dx$.\work{2.3cm}
  \item Find $\displaystyle\int\dfrac{1}{2x+1}\,dx$.\work{2.3cm}
\end{enumerate}

\level{Exam level}
\begin{enumerate}[leftmargin=*,itemsep=2pt,start=5]
  \item Find $\displaystyle\int x\sqrt{x^2+9}\,dx$.\work{2.8cm}
  \item Evaluate $\displaystyle\int_{0}^{\pi/2}\sin^2x\cos x\,dx$.\work{2.8cm}
  \item Evaluate $\displaystyle\int_{0}^{1}x^2\left(x^3+1\right)^3\,dx$.\work{2.8cm}
  \item Find $\displaystyle\int\cot x\,dx$ by writing $\cot x=\dfrac{\cos x}{\sin x}$.\work{2.8cm}
\end{enumerate}

\begin{frbox}{Question 9 \textperiodcentered{} Free-response practice}
Answer each part. No calculator is permitted.
\begin{enumerate}[label=(\alph*),leftmargin=*,itemsep=2pt]
  \item Evaluate $\displaystyle\int_{0}^{2}\dfrac{x}{x^2+4}\,dx$.\work{2.6cm}
  \item Evaluate $\displaystyle\int_{0}^{\pi}\sin x\,e^{\cos x}\,dx$.\work{2.6cm}
  \item If $\displaystyle\int_{0}^{8}f(x)\,dx=20$, find $\displaystyle\int_{0}^{2}f(4x)\,dx$.\work{2.6cm}
\end{enumerate}
\end{frbox}

\newpage
\begin{solbox}
\begin{enumerate}[leftmargin=*,itemsep=4pt]
  \item $u=x^2+1$: $\dfrac{\left(x^2+1\right)^4}{4}+C$.
  \item $\dfrac15e^{5x}+C$.
  \item $-\dfrac13\cos(3x)+C$.
  \item $\dfrac12\ln|2x+1|+C$.
  \item $u=x^2+9$: $\dfrac12\cdot\dfrac23u^{3/2}=\dfrac13\left(x^2+9\right)^{3/2}+C$.
  \item $u=\sin x$ from 0 to 1: $\Big[\dfrac{u^3}{3}\Big]_0^1=\dfrac13$.
  \item $u=x^3+1$ from 1 to 2: $\dfrac13\Big[\dfrac{u^4}{4}\Big]_1^2=\dfrac{15}{12}=\dfrac54$.
  \item $u=\sin x$, $du=\cos x\,dx$: $\ln|\sin x|+C$.
  \item \begin{enumerate}[label=(\alph*),leftmargin=*,itemsep=3pt]
    \item $u=x^2+4$ from 4 to 8: $\dfrac12\Big[\ln u\Big]_4^8=\dfrac12\ln2$.
    \item $u=\cos x$ from 1 to $-1$, $du=-\sin x\,dx$: $\displaystyle\int_{-1}^{1}e^{u}\,du=e-\dfrac1e$.
    \item $u=4x$, $dx=\dfrac14du$, limits 0 to 8: $\dfrac14(20)=5$.
  \end{enumerate}
\end{enumerate}
\end{solbox}

\end{document}
